% =====================================================================
%  符号说明(独立章节,编号随正文顺序)
%  由队长(Claude)统一维护,全队以此为准。
%  建模手写公式、代码手写代码、论文手排版都用同一套符号。
%  新符号出现时,更新此表并推送,大家在群里同步。
%
%  排版说明(国赛格式规范):单列表,每行一组「符号 含义 单位」,
%  列宽 符号 3.2cm / 含义 10.0cm / 单位 1.3cm,
%  \small 字号 + \tabcolsep=2.5pt;[H] 钉位,表间连排(不 \clearpage,避免大段空白)。
%  问题二、问题三的新符号在本文件末尾按新表追加。
% =====================================================================
\section{符号说明}

% ---------- 表一:问题一(残损存在性分类) ----------
\begin{table}[H]
\centering
\small
\renewcommand{\arraystretch}{1.0}
\setlength{\tabcolsep}{2.5pt}
\caption{主要符号说明(一):问题一残损存在性分类}
\begin{tabular}{p{3.2cm}p{10.0cm}p{1.3cm}}
\toprule
符号 & 含义 & 单位 \\
\midrule
$I$ & 输入的集装箱外表面图像($640\times640$ RGB) & — \\
$f(I)$ & 残损存在性分类函数,取值 $\{0,1\}$ & — \\
$D(\cdot)$ & 残损检测模型(问题二建立,问题一复用) & — \\
$\mathcal{B}$ & 检测器输出的候选框集合 & — \\
$K$ & 候选框总数($\mathcal{B}=\{(b_i,c_i,s_i)\}_{i=1}^{K}$) & — \\
$b_i$, $c_i$, $s_i$ & 第 $i$ 个检测框的位置 $(x_i,y_i,w_i,h_i)$、类别(0 凹陷/1 破洞/2 锈蚀)与置信度 & — \\
$\tau$ & 判定“存在残损”的置信度阈值 & — \\
$N(I;\tau)$ & 置信度不低于 $\tau$ 的检测框个数 & — \\
$S(I)$ & 图像级最高检测置信度($\max_i s_i$,无框记 0) & — \\
$Acc$, $P$, $R$ & 分类准确率、查准率、召回率 & — \\
$F_1$, $AUC$ & $F_1$ 值、ROC 曲线下面积 & — \\
\bottomrule
\end{tabular}
\end{table}

% ---------- 表二:问题二(检测与分割指标) ----------
\begin{table}[H]
\centering
\small
\renewcommand{\arraystretch}{1.0}
\setlength{\tabcolsep}{2.5pt}
\caption{主要符号说明(二):问题二检测与分割指标}
\begin{tabular}{p{3.2cm}p{10.0cm}p{1.3cm}}
\toprule
符号 & 含义 & 单位 \\
\midrule
$B^{gt}$ & 真实(ground-truth)边界框 & — \\
$\mathrm{IoU}$ & 交并比 & — \\
$P$, $R$ & 查准率、召回率 & — \\
$AP$ & 平均精度(PR 曲线下面积) & — \\
$mAP$ & 均值平均精度 & — \\
$mAP@0.5$ & 交并比阈值 0.5 下的均值平均精度 & — \\
$mAP@0.5:0.95$ & 交并比阈值从 0.5 到 0.95 平均的均值平均精度 & — \\
$\hat{M}_c$, $M_c$ & 预测掩码与真实掩码(类别 $c$) & — \\
$\hat{m}_i$, $m_i$ & 掩码像素预测值与真值 & — \\
$\mathrm{Dice}$ & Dice 系数(分割) & — \\
$\mathrm{mIoU}$ & 平均交并比(分割) & — \\
\bottomrule
\end{tabular}
\end{table}

% ---------- 表三:问题二(网络损失与后处理) ----------
\begin{table}[H]
\centering
\small
\renewcommand{\arraystretch}{1.0}
\setlength{\tabcolsep}{2.5pt}
\caption{主要符号说明(三):问题二网络损失与后处理}
\begin{tabular}{p{3.2cm}p{10.0cm}p{1.3cm}}
\toprule
符号 & 含义 & 单位 \\
\midrule
$\mathcal{L}_{cls}$ & 分类损失 & — \\
$\mathcal{L}_{box}$ & 边框回归损失 & — \\
$\mathcal{L}_{seg}$ & 分割损失 & — \\
$\mathcal{L}$ & 总损失 & — \\
$\lambda_{cls}$,$\lambda_{box}$,$\lambda_{seg}$ & 各损失项权重 & — \\
$\tau_{nms}$ & 非极大值抑制(NMS)的 IoU 阈值 & — \\
$sev(b_i)$ & 检测框严重度代理($=w_i\,h_i$,面积越大越严重) & — \\
$P_2$, $P_3$, $P_4$, $P_5$ & 特征金字塔不同尺度的特征层($P_2$ 为小目标增强层) & — \\
\bottomrule
\end{tabular}
\end{table}

% ---------- 表四:问题三(模型评估) ----------
\begin{table}[H]
\centering
\small
\renewcommand{\arraystretch}{1.0}
\setlength{\tabcolsep}{2.5pt}
\caption{主要符号说明(四):问题三模型评估}
\begin{tabular}{p{3.2cm}p{10.0cm}p{1.3cm}}
\toprule
符号 & 含义 & 单位 \\
\midrule
$TPR$, $FPR$ & 真阳率、假阳率 & — \\
$AP_s$, $AP_m$, $AP_l$ & 小、中、大目标平均精度 & — \\
$y_j^{clean}$, $y_j^{pert}$ & 干净/扰动条件下第 $j$ 项指标取值 & — \\
$\Delta_j$ & 第 $j$ 项指标的性能衰减率 & \% \\
$t_{inf}$ & 单张图像平均推理耗时 & s \\
$FPS$ & 推理帧率 & 帧/s \\
$A^k$ & 第 $k$ 个特征图 & — \\
$y^c$ & 类别 $c$ 的预测得分 & — \\
$\alpha_k^c$ & 第 $k$ 个特征图对类别 $c$ 的权重 & — \\
$G^c$ & 类别 $c$ 的 Grad-CAM 激活图 & — \\
\bottomrule
\end{tabular}
\end{table}

% 说明:后续新增符号时:告知队长 → 队长更新本表 → 推送,全队同步。
